\section{The Census Nobody Has}

The Church has a formal procedure for declaring that a particular person is in heaven. It is called canonization, it is used several times a year, it has been used many thousands of times, and its effect is precisely a statement about an individual's actual condition after death. There is no corresponding procedure in the other direction. There never has been. No pope, no council, and no rite has ever declared of any human being who has ever lived that he is in hell.

That is a negative claim about a very large corpus, and it should be stated with its bounds. This essay did not read every magisterial act; what it did was read the eight acts of the Catechism's own dossier, the paragraphs of the Catechism on hell and on mortal sin, the relevant canons of Trent, and the most recent papal statement on the subject. None of them names anyone. And the claim happens to be made in that most recent statement, by the pope, as part of the doctrine: ``Damnation remains a real possibility, but it is not granted to us, without special divine revelation, to know which human beings are effectively involved in it.''\endnote{John Paul II, general audience of 28 July 1999, n.~4; official Vatican English web text, checked 2026-07-25. The sentence occurs immediately after the paragraph on the fallen angels, which is deliberate and is discussed below. A general audience is not a definition; but the sentence is a report of what the Church holds herself able to say, made by the person in the best position to know, and no contrary act is known to this essay.}

The rest of this section is about why the census is empty --- which turns out not to be a policy of discretion but a consequence of the doctrine itself.

\subsection{The one place the question is asked}

Somewhere on the road to Jerusalem a man asked Jesus the question that everyone since has wanted answered, and Luke records the exchange in two verses:

\begin{studyframe}
And a certain man said to him: Lord, are they few that are saved? But he said to them: Strive to enter by the narrow gate: for many, I say to you, shall seek to enter and shall not be able.\endnote{Luke 13:23--24, Douay--Rheims, read 2026-07-25 in the tracked verse-text derivative with the whole of 13:22--30. Greek from the tracked Robinson--Pierpont CSV for Luke: \latin{Eipen de tis aut\=oi, Kyrie, ei oligoi hoi s\=ozomenoi? Ho de eipen pros autous, Ag\=onizesthe eiselthein dia t\=es sten\=es pyl\=es\ldots{}}}
\end{studyframe}

Read the grammar. One man, singular, asks a question in the third person about other people: are the saved few? The answer is addressed to a plural audience and is an imperative in the second person: \emph{you}, strive. \latin{Ag\=onizesthe} is the verb of the wrestling floor. The question asked for a number and got a command; the question was about them and the answer was about you.

That is the only place in the Gospels where the population of the saved is put to Christ as a direct question, and the shape of the reply is the shape of the whole doctrine. It is not that the question was rebuked --- there is no rebuke in the text. It is that it was answered on a different axis, and that the axis it was answered on is the only one anybody can act on.

The other numerical texts are of the same kind, and it is worth reading them for their grammatical mood rather than their arithmetic. ``Enter ye in at the narrow gate: for wide is the gate, and broad is the way that leadeth to destruction, and many there are who go in thereat. How narrow is the gate, and strait is the way that leadeth to life: and few there are that find it!''\endnote{Matthew 7:13--14, Douay--Rheims at its locus, read 2026-07-25. The passage stands in the Sermon on the Mount and opens with an imperative, \latin{Eiselthete}; the ``many'' and ``few'' clauses are subordinate to it. The Catechism quotes exactly this text at CCC 1036 and introduces it as ``an urgent call to conversion,'' which is a statement about its function.} The sentence does contain ``many'' and ``few,'' and an honest reading does not pretend otherwise; the way to life is described as the one fewer people find. But the sentence is not a report, it is a command with a reason attached, and it is addressed to the hearer about the hearer. Neither text has ever been read by the magisterium as a demographic statement, and the Catechism, quoting the second of them in the paragraph immediately after its statement of the doctrine, glosses it as ``an urgent call to conversion.''

\subsection{Judas, and the limit of what can be said}

If the Church were ever going to name anyone, the case is made for her. Of one man the Lord said: ``But woe to that man by whom the Son of man shall be betrayed. It were better for him, if that man had not been born.''\endnote{Matthew 26:24, Douay--Rheims at its locus, read 2026-07-25. The parallel at Mark 14:21 was also read.} Of one man he said, in the priestly prayer: ``none of them is lost, but the son of perdition.''\endnote{John 17:12, Douay--Rheims at its locus, read 2026-07-25. The phrase rendered ``son of perdition'' is a Semitic idiom of belonging; the Greek was not consulted for this verse in a registered edition, and no argument here rests on its wording.} Of one man the apostles said, choosing his successor, that he ``hath by transgression fallen, that he might go to his own place.''\endnote{Acts 1:25, Douay--Rheims at its locus, read 2026-07-25.} Three texts, of unusual severity, about a named individual.

The tradition has overwhelmingly read them as they appear. Augustine, Aquinas, and the general run of Latin theology treat Judas as lost, and a reader who finds that reading natural is in very large company. What has never happened is the next step. No council has defined it, no pope has declared it, and the Catechism, which quotes both Matthew 26 and the Gospel warnings about hell, makes no statement about Judas at all. On the one human being about whom Scripture speaks in this way, the Church's formal teaching is silent.

Two conclusions are available from that silence and only one of them is warranted. The unwarranted one is that the Church privately doubts the doctrine and is hedging. The warranted one is that the act of declaring an individual damned is not an act she has ever understood herself to be able to perform --- which is a claim not about Judas but about the kind of knowledge that would be required.

\subsection{Why the knowledge is not available}

The Church's own analysis of mortal sin makes the private judgment structurally impossible, and this is the load-bearing point of the whole section. Three conditions must be met together: ``Mortal sin is sin whose object is grave matter and which is also committed with full knowledge and deliberate consent.'' The first condition is public --- grave matter is specified by the commandments, and whether an act falls under one of them is a question anybody can examine. The second and third are not:

\begin{studyframe}
Unintentional ignorance can diminish or even remove the imputability of a grave offense\ldots{} the promptings of feelings and passions can also diminish the voluntary and free character of the offense, as can external pressures or pathological disorders.\endnote{CCC 1857 and 1860, official Vatican English web text, checked 2026-07-25 together with CCC 1854--1864. The paragraph continues: ``But no one is deemed to be ignorant of the principles of the moral law, which are written in the conscience of every man,'' and ``Sin committed through malice, by deliberate choice of evil, is the gravest.'' The Church's analysis is not an amnesty; it is a statement of what would have to be known.}
\end{studyframe}

Knowledge and consent are interior. Whether a particular man knew what he was doing, and whether the assent he gave was free rather than driven by compulsion, terror, formation, or illness, is not visible from outside and is frequently not visible from inside. And even if it were, there remains the last instant, which is invisible to everyone: the doctrine is about dying unrepentant, and the Church's phrase is always \emph{to the end}. John Paul II puts the two together: hell is the state of those who ``definitively reject the Father's mercy, \emph{even at the last moment of their life}.''\endnote{John Paul II, general audience of 28 July 1999, n.~1; same witness and date. The emphasis is added here. Compare Paul VI's \emph{Credo} n.~12, ``those who have refused them to the end,'' quoted in the section on the acts.}

The Catechism draws the practical conclusion in the same paragraph in which it states that mortal sin causes ``the eternal death of hell'': ``However, although we can judge that an act is in itself a grave offense, we must entrust judgment of persons to the justice and mercy of God.''\endnote{CCC 1861, official Vatican English web text, checked 2026-07-25. The full sentence preceding it is: ``If it is not redeemed by repentance and God's forgiveness, it causes exclusion from Christ's kingdom and the eternal death of hell, for our freedom has the power to make choices for ever, with no turning back.'' The two halves of the paragraph belong together and are regularly quoted apart.} That is not a counsel of tact. It is a statement about the location of a boundary: acts can be judged, persons cannot, and the doctrine of hell concerns persons.

The same limit is imposed even on self-knowledge, which is the point at which the Church's discipline becomes genuinely surprising. Trent, having taught that the justified should place their firmest hope in God's help, adds that concerning the gift of perseverance ``let no one promise himself anything as certain with absolute certainty'' --- \latin{nemo sibi certi aliquid absoluta certitudine polliceatur} --- and directs the reader instead to Paul: let him that thinketh himself to stand take heed lest he fall, and work out his salvation with fear and trembling.\endnote{Council of Trent, session 6, chapter 13 on the gift of perseverance; DS 1541 (DB 806), Latin read 2026-07-25 at patristica.net, the chapter read whole. The corresponding canon is canon 16 of the same session. Trent is not teaching despair: the same chapter says that all must place and repose a most firm hope in God's help, and that God, unless men fail his grace, will complete the good work he has begun.} If a man cannot have absolute certainty about the one case he can see from the inside, the case for which he has continuous first-person evidence and unlimited access, then the epistemic situation with respect to anyone else is not merely difficult. It is closed.

\subsection{One class about which the Church does speak}

Precision requires an exception, and it is instructive. Everything above concerns human beings. About the fallen angels the Church's language is not conditional at all.

Canon 9 of 543 anathematizes the proposition that the punishment ``of demons and of impious men'' will end --- naming the demons first. \emph{Firmiter} says the wicked receive perpetual punishment \latin{cum diabolo}, with the devil. Matthew 25:41 has the fire ``prepared for the devil and his angels,'' in the perfect participle: it exists, and it exists for them. The Apocalypse says of the devil, the beast, and the false prophet that ``they shall be tormented day and night for ever and ever.'' John Paul II states it as doctrine in the same address that refuses the human census: ``Christian faith teaches that in taking the risk of saying `yes' or `no', which marks the human creature's freedom, some have already said no. They are the spiritual creatures that rebelled against God's love and are called demons.''\endnote{John Paul II, general audience of 28 July 1999, n.~4, citing the Fourth Lateran Council, DS 800--801; same witness and date. The paragraph ends: ``What happened to them is a warning to us.'' Lateran IV's chapter, read above, states that the devil and the other demons were created good in nature and became evil of themselves.}

So the Church does hold that hell is occupied. What she has never said is that any \emph{human being} is in it. The distinction is not arbitrary: the angelic choice, on the traditional analysis, was made once, with full knowledge, by an intellect that does not deliberate as we do --- which is exactly the case in which the three conditions of a mortal sin are transparently met, and exactly the case that is furthest from every human case anyone could examine. The one class about which the doctrine speaks in the indicative is the one class whose interior act is not opaque.

\subsection{What follows for speech}

The consequences for how the doctrine may be spoken are strict and cut in a direction that both parties in the modern quarrel dislike.

Nobody may say of any person, living or dead, that he is damned. Not of an enemy, not of a public criminal, not of a suicide, not of an apostate, not of anyone. This is not sentiment; it is the direct consequence of a rule the Catechism states and Trent enforces even against the speaker's knowledge of himself. Preaching that consigns named individuals, or identifiable classes of individuals, to hell is not rigorous orthodoxy. It is a claim to information the Church says nobody has.

And nobody may say that the doctrine has no application. The absence of a census is not evidence of an empty hell; it is the absence of a census. The definitions are in the indicative --- souls departing in mortal sin descend immediately --- and the conditional under which they operate is one that the Church treats as really satisfiable, or the warning would be a pretence and the Roman Canon's petition to be saved from final damnation would be a formality.

There is a third consequence, which cuts against a modern reflex and is worth stating because it is definitively taught. The fear of hell is not a disreputable motive. Trent anathematized the proposition that the sorrow ``which is prepared by the examination, recollection, and detestation of sins,'' by which a man weighs ``the gravity, multitude, and foulness of his sins, the loss of eternal beatitude, and the incurring of eternal damnation,'' together with a purpose of amendment, ``is not a true and useful sorrow, nor prepares for grace, but makes a man a hypocrite and a greater sinner.''\endnote{Council of Trent, session 14 (25 November 1551), canon 5 on the sacrament of penance; DS 1705 (DB 915), Latin read 2026-07-25 at patristica.net: \latin{\ldots{}amissionem aeternae beatitudinis, et aeternae damnationis incursum, cum proposito melioris vitae, non esse verum et utilem dolorem, nec praeparare ad gratiam, sed facere hominem hypocritam et magis peccatorem: an.~s.} The corresponding doctrinal chapter is DS 1677--1678. The canon does not teach that fear is the best motive, and the tradition is unanimous that contrition from charity is higher; it teaches that fear of damnation is a real motion of grace and not a manufacture of hypocrites.} The charge that the doctrine produces frightened hypocrites rather than converts is exactly the charge Justin Martyr was answering in the second century, and it is one of the very few objections in this whole subject that the Church has met with a formal anathema.
